\documentclass[11pt,a4paper]{article}

%used packages
\usepackage{graphicx}
\usepackage[dutch]{babel} 
\usepackage{scrextend}
\usepackage{enumitem}
\usepackage{listings}
\usepackage{float}
\usepackage{titlesec}
\usepackage[table,xcdraw]{xcolor}
\usepackage{fancyhdr}
\usepackage{lastpage}
\usepackage{xcolor}
\usepackage{listings}


%Document settings
\raggedbottom
\graphicspath{ {images/} }
\setlength{\parindent}{0em}
\setlength{\parskip}{1em}

%Custom commando's
\newcommand\litem[1]{\item{\bfseries#1.\space}}
\newcommand\tabelkleur{\rowcolor[HTML]{FFCC67}}


%Custum Variables
\def\auteureen{Roy Buitenhuis, 0895833}
\def\auteurtwee{Tim van Broekhoven, 0893122}
\def\titel{Verslag TDS01} 
\def\datum{\today}
\def\versie{0.1}
\def\status{concept}
\def\subtitel{Plan van aanpak}
\def\bedrijf{}
\def\tabelkleur{FFCC67}

%Settings for Footer and Header
\pagestyle{fancy}
\fancyhf{}
\rhead{\subtitel}
\lhead{\titel}
\lfoot{\datum}
\rfoot{Pagina \thepage \hspace{1pt} /  \pageref{LastPage}}

%Settings for document spacing
\titlespacing{\section}{0pt}{*0}{*0}
\titlespacing{\subsection}{0pt}{*0}{*0}
\titlespacing{\subsubsection}{0pt}{*0}{*0}

%Settings for code sinppeds
\lstset{basicstyle=\ttfamily,
	showstringspaces=false,
	commentstyle=\color{red},
	keywordstyle=\color{blue}
}

\begin{document}
	
	\begin{titlepage}
		
		\centering
		\includegraphics[width=0.80\textwidth]{image1}\par\vspace{1cm}
		{\huge\bfseries \titel \par}
		
		\vspace{1cm}
		{\Large\itshape \auteureen \par}
		{\Large\itshape \auteurtwee \par}
		\vspace{1cm}
		{\Large\itshape versie \versie\par}
				
		\vfill
		Vak:\par
		TDS01
		
		\vfill
		{\large \datum \par}
	\end{titlepage}

	\section{Samenvatting}
	
	\clearpage
	
	\tableofcontents
	
	\clearpage
	
	\listoffigures
	
	\clearpage
	\listoftables
	
	\clearpage
	
	\section{Versiehistorie}
	\begin{table}[H]
		\centering
		\label{Versiehistorie}
		\begin{tabular}{|p{1cm}|p{2cm}|p{6cm}|p{2cm}|}
			\hline
			\rowcolor[HTML]{FFCC67}
			\textbf{Versie} & \textbf{Datum} & \textbf{Wijzigingen} & \textbf{Auteur} \\ \hline
			0.1    & 07-09-2017 & Template    & Thijs \\ \hline
			0.2	   & 19-09-2017 & Oplevering eerste versie  & Groep \\ \hline
			&       &             &        \\ \hline
		\end{tabular}
		\caption {Versiehistorie} \label{tab:title} 
	\end{table}	


	\section{Introductie}
		\subsection{Het vak}
		TDS02 is een van de vakken die tijdens de minor 'Embedded Systems' wordt gegeven. Het vak bestaat voornamelijk uit practicum assignments die de studenten in groepjes van twee dienen te voltooien. In de eerste weken begint de les met een uitleg van de docent over de theorie achter deze assignments, die het doel hebben om de studenten te trainen in digitale signaalbewerking. Om de assignments te voltooien dienen de studenten gebruik te maken van de 'C5505 eZdsp Development Tool' van Texas Instruments.
		
		\subsection{Introductie DSP}
		Volgens Analog Devices \cite{analog} is een DSP een processor die een  	gedigitaliseerd signaal als geluid, video, temperatuur of positie op wiskundige wijze manipuleerd. Analog Devices ligt toe dat een DSP wordt ontworpen om berekeningen als "optellen", "aftrekken", "vermenigvuldigen" en "delen" in korte tijd te kunnen voltooien. Het doel van deze berekeningen is om met de informatie van het ingangssignaal een uitgangssignaal te produceren die bruikbaar is voor een bepaalde toepassing. 
		
		\subsection{Digitale filter}
		Een digitale filter bewerktt een digitaal en tijd-discrete signaal om bepaalde frequenties uit dat signaal te verwijderen.
		
		Bron: http://www.dspguide.com/ch14/1.htm
		boek ook in pdf  			
		
		\subsection{De opdracht}		
		
				
	
	\section{Ontwerp en realisatie FIR filter}
	
		\subsection{Requirements FIR filter}
	
		\subsection{Coefficienten bepalen}
	
		\begin{enumerate}[label=\emph{\alph*)}]
			\litem{item} Describe the settings you 					used in MATLAB’s FDAtool and explain the 					choices you made.
		\end{enumerate}
	
		\subsection{Matlab}
	
		\subsection{Code}

		\begin{lstlisting}
			%code
		\end{lstlisting}
		
		\subsection{Het resultaat}

	
	\section{Ontwerp en realisatie IIR filter}
		
		\subsection{Requirements IIR filter}
	
		\subsection{Coefficienten bepalen}
	
		\begin{enumerate}[label=\emph{\alph*)}]
			\litem{item} Describe the settings you 					used in MATLAB’s FDAtool and explain the 					choices you made.
		\end{enumerate}
	
		\subsection{Matlab}
	
		\subsection{Code}

		\begin{lstlisting}
			%code
		\end{lstlisting}
		
		\subsection{Het resultaat}	
	
	\section{optimizatie}
	

	\section{Conclusie en aanbevelingen}	

	
	\bibliographystyle{plain}
	\bibliography{bibligraphy}
		
	\begin{thebibliography}{9}

	\bibitem{analog}
  		
  		\textit{A Beginner's Guide to Digital Signal Processing (DSP)},
  		Analog Devices,
  		URL: http://www.analog.com/en/design-center/landing-pages/001/beginners-guide-to-dsp.html 

	
	\bibitem{}
  		
  		\textit{},
  		Addison Wesley, Massachusetts,
  		2nd edition,
  		1994.
  		
	\bibitem{}
  		
  		\textit{},
  		Addison Wesley, Massachusetts,
  		2nd edition,
  		1994.  		
		
	\end{thebibliography}

\end{document}